\section{Introduction: Why Pattern Realism Needs a Stricter
  Criterion}\label{introduction-why-pattern-realism-needs-a-stricter-criterion}

The special sciences routinely succeed by tracking macro-descriptions
rather than exhaustively tracking microstates. When that success is
stable under perturbation and supports reliable intervention, the
natural reading is realist: the macro-description is tracking
something in the world. But no existing philosophical framework
provides a sharp dynamical criterion for when this realist reading is warranted
and when it is not. Dennett's real-pattern realism comes closest.
His criterion is explicit: ``A pattern exists in some data, and is
real, if there is a description of the data that is more efficient
than the bit map, whether or not anyone can concoct it''
\citep[p. 34]{dennett1991}. That captures something correct about the
structure of explanation. The difficulty is that compression alone is
too permissive. A coding strategy can be general and
compressive while still grouping together microstates whose onward
transition profiles diverge. Such constructions compress, yet they
require persistent within-class bookkeeping to maintain forecast
quality, because the macro-label hides transition-relevant
heterogeneity. This residual permissiveness is the target of the
present paper.

The worry has a well-established lineage. Haugeland pressed the
question of when pattern detection crosses from descriptive
convenience into ontological contact \citep{haugeland1998},
and the Putnam-Fodor autonomy tradition established that higher-level
kinds can be explanatorily indispensable \citep{putnam1967,fodor1974},
but left implicit the dynamical condition that separates genuine higher-level kinds from merely disjunctive aggregates. More recent
work on idealization continues to press the gap between descriptive
utility and epistemic warrant \citep[chap.~2]{elgin2017}. The present
paper takes that composite critique seriously. What is missing, on the present view, is an explicit test of whether macro-descriptions carry
autonomous transition structure, which further compression cannot
supply.

A contrast makes the gap concrete. Lewis observed that we have no
principled reason to deny existence to the mereological sum of ``the
right half of my left shoe plus the Moon plus the sum of all Her
Majesty's ear-rings,'' however sensible it is to ignore such things in
ordinary thought \citep[p. 213]{lewis1986}. That is correct as far as it
goes. But consider the difference between that composite and a
hurricane. Knowing a hurricane's pressure organization and rotational
structure tells you what it will do next without tracking every
molecule. The shoe-moon composite carries no comparable transition
structure; predicting its future requires independently tracking each
component. On Lewis's view both exist; what the hurricane has and
the composite lacks is autonomous macro-dynamics. This paper argues that closure of transition structure
is the criterion that marks that difference.

The thesis is that real-pattern realism needs an explicit closure
condition to avoid permissive drift. For a fixed regime, horizon, and admissible intervention class, a
candidate macro-object qualifies when macrostate information is
sufficient for macro-transitions, so that within-class micro-differences
make no difference to what happens next at the macro level.

Throughout this paper, \emph{gerrymandered} means dynamically
incoherent: a partition that fails transition autonomy by requiring
persistent within-class micro-bookkeeping. A partition that is merely
intuitively disjunctive or aesthetically odd does not count as
gerrymandered in this sense.

The thesis tightens existing pattern realism and keeps its insights
about compression and projectibility. What it adds is a discriminating condition that excludes
gerrymandered candidates by transition structure rather than by
intuitive naturalness. The formal benchmark for this condition is strong
lumpability in exact Markov settings. The non-ideal extension is graded
closure via leakiness and convergent diagnostics under fixed
constraints.

The argument is philosophical, not a methods paper in disguise. The paper provides a criterion and a
disciplined protocol for applying it; it does not claim to deliver a
universal estimation recipe.
Estimator selection, finite-sample behavior, and domain-specific
implementation remain downstream methodological tasks.

\subsection{Novelty and Positioning}\label{novelty-and-positioning}

The novelty claim has three parts.

\begin{enumerate}
	\item
	      Beyond Dennett alone: compression realism is retained, but
	      permissiveness is reduced by an explicit anti-gerrymandering
	      condition.
	\item
	      Beyond formal closure results alone: recent work distinguishes
	      informational, causal, and computational closure as formal
	      diagnostics \citep{rosas2024}, but does not say when those
	      diagnostics warrant ontological commitment or when commitment
	      should be withdrawn. The present paper supplies that missing
	      philosophical layer: an objecthood criterion with explicit
	      intervention indexing, graded verdicts, and downgrade conditions.
	\item
	      Beyond pure interventionist pragmatism: admissibility is fixed
	      upstream and physically constrained, so verdicts are not back-fit to
	      analyst preference.
\end{enumerate}

This is a conditional metaphysical proposal. Under structural realist
and interventionist commitments, closure under admissible conditions is
sufficient for macro-objecthood in regime. Readers who withhold those
commitments can still accept a narrower conclusion: closure is at least
a necessary anti-gerrymandering constraint on serious macro-ontology claims.

With Dennett, the paper
keeps compression realism while rejecting permissiveness \citep{dennett1991}. With Ladyman and Ross, it keeps structural and projectibility
ambitions while insisting on an explicit transition criterion for
objecthood verdicts \citep{ladyman2007}. With Rosas et al., the paper uses their formal closure diagnostics as
measurement machinery and supplies the philosophical
criterion that determines what those measurements mean for ontology
\citep{rosas2024}. With causal-emergence work, it treats macro-level
gains as corroborating diagnostics rather than as a standalone ontology
test \citep{hoel2013,dewhurst2021}. With Kim-style exclusion pressure, it defends
non-redundancy at explanatory and control levels without denying
microphysical implementation \citep{kim1998}. One consequence of
intervention indexing is that the present criterion can diverge from
purely observational diagnostics: a macro-variable can achieve
predictive sufficiency by tracking a common cause rather than by
instantiating autonomous transition structure, passing observational
tests while failing under admissible intervention
(\S\ref{what-the-formal-machinery-alone-does-not-provide}).

Three further traditions bear on the view. Simon's
account of hierarchy and near-decomposability explains why complex
systems can display relatively autonomous short-run subsystem dynamics
alongside looser longer-run interdependence \citep{simon1962,ando1961}.
That is a structural precursor to the regime- and horizon-indexed
treatment of closure defended here. Simon asks when aggregation is
dynamically serviceable; he leaves open when such serviceability
warrants ontological commitment, which is the question pursued here.

Wimsatt's analysis of aggregativity marks a second antecedent. He argues
that fully aggregative system properties are exceptional because most
higher-level properties depend on organization and therefore fail
invariance conditions such as unrestricted rearrangement,
reaggregation, and intersubstitution of parts \citep{wimsatt2000}. That
diagnosis fits the paper's anti-gerrymandering aim, but
non-aggregativity alone does not show that a candidate macro-description
carries self-contained transition structure. Closure adds that test.

The mechanisms tradition provides a third point of contact. Machamer,
Darden, and Craver characterize mechanisms in terms of organized
entities and activities productive of regular change, while Craver's
later account emphasizes that explanatory levels are local to
mechanisms rather than pieces of one global hierarchy
\citep{machamer2000,craver2007}. The present paper is continuous with
that tradition, but adds a criterion for when a candidate variable or
coarse-graining has earned a place in a mechanism-level ontology at all.

A related line of work traces the cognitive origins of real-pattern
commitments, linking predictive coding and free-energy minimization to
how bounded agents build and revise their representational schemes
\citep{gladziejewski2025}. That work helps explain why agents converge on
some macro-descriptions and abandon others, which strengthens the
anti-arbitrariness side of pattern realism. But representational success
can remain task-relative unless a further test determines when a
candidate grouping tracks dynamical structure in the world rather than
merely serving local inferential needs. Closure supplies that test.

A stronger reading of Dennett requires a reply. One can
interpret real-pattern realism as already requiring projectibility and
cross-context robustness, not merely compression in a single dataset.
On that steelman, Dennett's criterion already excludes fragile,
narrowly-tuned codings. Even so, projectibility constrains which
predictions are worth tracking; it does not test whether the
macro-description carries autonomous transition structure. A partition
can be projectible across contexts and still require persistent
within-class micro-bookkeeping to sustain its forecasts, because the
macro-label groups together microstates whose onward profiles diverge
once perturbation or horizon extension is applied. Closure catches
that residual failure mode, because in addition to asking whether a
pattern projects, it asks whether onward macro-transitions are
self-contained at the declared grain. Even on the most charitable
reading of Dennett, then, projectibility can still tolerate successful
forecasting that depends on recurring micro-level repair, and closure
rules such forecasting out as a basis for objecthood.

\subsection{Scope and Non-Claims}\label{scope-and-non-claims}

Closure can be stipulated or induced. In stipulated closure, autonomy is fixed by constitutive
rules, as in chess or formal arithmetic: once the rules are set,
transitions and consequences follow by construction. In induced
closure, autonomy is discovered through dynamical screening-off in the
physical world. This paper is concerned exclusively with induced
closure in spatiotemporal systems. It does not offer a complete metaphysics of levels, and it does not settle
all questions about abstract objects or full mereology. Where those
debates arise, they are handled only to protect the central claim from
predictable misreadings. Still, the criterion is stated over
state-transition structure: spatiotemporal dynamics are the paradigm
case, not the only intelligible format in which closure questions can be
posed.

The claim that closure is sufficient for macro-objecthood covers the
question of whether there is
an autonomous macro-level unit at all: a pattern that behaves as a
coherent whole at its grain, with its own transition structure. It does
not cover everything ordinary talk about objects bundles together, such
as what kind of thing the unit is, how it came about, what it is for,
how it is named, or how it is traced through historical and social
change. Those questions belong to other accounts, and several current
ones fit alongside closure without competing with it. Contextual
emergence treats higher-level descriptions as licensed by contingent
contexts that lower-level description does not itself supply
\citep{bishopatmanspacher2006}; the regime-indexing defended here is a
version of that idea. Process ontology treats organisms and other
entities as stable organizations of ongoing change \citep{dupre2018},
which is also how closure treats them: what persists is a transition
organization, not a fixed collection of matter. Structurally oriented
hylomorphism fits to the extent that ``form'' is read as the organizing
structure a substrate realizes \citep{jaworski2016}. Accounts that tie
objecthood to physical binding describe one important way closure is
produced in some physical domains. Functional and historical accounts of
artifacts explain why something is classified as a hammer or a chair
\citep{thomasson2007}. These accounts explain how closed structure is
realized, stabilized, classified, and individuated. Closure says whether
the unit they describe is there to be realized, stabilized, classified,
and individuated in the first place, and it says so in terms of
transition structure.

Methodologically, the paper does not depend on the strongest
anti-analytic rhetoric sometimes associated with naturalized
metaphysics. The claim is narrower: closure is a usable philosophical
constraint that can be motivated by dynamical and coarse-graining
practice, whether or not one accepts every meta-metaphysical commitment
in the broader \emph{Every Thing Must Go} (ETMG) program \citep{ladyman2007}.

Section 2 states the criterion and admissibility
discipline. Section 3 gives the exact benchmark. Section 4 extends the
criterion to non-ideal settings. Section 5 addresses the strongest
objections. Section 6 summarizes gains and limits.

\subsection{Three Nearby Positions}\label{three-nearby-positions}

Three nearby stances are often run together.

\begin{enumerate}
	\item
	      Compression-prediction realism: successful compression and forecasting
	      are treated as sufficient for realist commitment.
	\item
	      Pure pragmatism: model choice is guided by utility alone, with no
	      ontological consequence.
	\item
	      Closure realism: compression and prediction matter, but on the
	      present account objecthood requires transition autonomy relative to a fixed, declared regime, horizon, and admissible intervention class.
\end{enumerate}

The third stance preserves the explanatory strengths of the first while
avoiding its permissiveness. It also avoids reducing ontology to
convenience, which is the pressure on the second.

\subsection{Contribution Map}\label{contribution-map}

The paper makes four linked contributions.

\begin{enumerate}
	\item
	      It gives a criterion-level tightening of real-pattern realism by
	      adding closure as a discriminating objecthood condition.
	\item
	      It provides a bridge from exact benchmark cases to non-ideal cases
	      without changing criterion content.
	\item
	      It offers a disciplined verdict structure with explicit downgrade
	      conditions, so realism claims remain answerable to failure.
	\item
	      It shows that closure supplies the dynamical condition the
	      Putnam-Fodor autonomy tradition needed but left implicit:
	      multiple realizability groups realizers extensionally, while
	      closure demands that those realizers produce the same
	      macro-transitions (\S\ref{autonomy-pluralism-and-distinctiveness}).
\end{enumerate}

The central payoff is selective commitment. The account can support
strong commitments where transition autonomy is robust, while avoiding
forced commitments where evidence is fragile or regime-sensitive.

\section{Closure and Admissibility}\label{closure-and-admissibility}

\subsection{Criterion in Plain
	Language}\label{criterion-in-plain-language}

Closure can be stated without heavy formalism. A candidate macro-object
passes when knowing its current macrostate is enough to determine
its macro-future over a specified horizon and intervention
class. If two microstates inside one macrostate produce different
macro-transitions, closure fails at that grain.

Static patterns fall under the same rule. In a static pattern,
micro-fluctuations are screened off rather than absent, so persistence
is itself a kind of macro-transition and falls within the closure
criterion. The criterion concerns transition sufficiency, not
descriptive fit alone. A representation can summarize trajectories
elegantly and still fail objecthood if it requires persistent
within-class micro-bookkeeping to preserve forecast quality.

The formal machinery evaluates partitions of state spaces, not physical
things directly, so it needs an explicit link to objects. On the
structural realist commitments defended below
(\S\ref{why-closure-carries-ontological-weight}), a macro-object just is
stable relational and causal organization at the relevant grain; the
partition identifies the object by identifying that organization.

Talk of ``tests'' and ``criteria'' in this paper is epistemic talk, and
it should not be read back into the metaphysics. Whether a partition
carries autonomous transition structure under a given regime, horizon,
and intervention class is a fact about the system. It holds or fails
whether or not anyone has checked. The regime is not up to us either.
It is picked out by physical features of the system, such as
timescales, energy scales, and available control channels, that can be
identified without knowing which candidate partition will close
(\S\ref{regime-individuation-and-circularity}). What investigators
contribute is a claim about where such regimes and closed partitions
lie, and that claim can be wrong. The procedural rules below, such as
fixing the regime, horizon, and admissible class before scoring, govern
the claim. They keep the evidence from being tuned to a favored answer,
and they play no part in making the facts. The indexing itself works
like the temperature and pressure conditions on the existence of ice:
it says which modal claim is being made, and it does not make the
verdict depend on anyone's interests.

\subsection{Formal Statement}\label{formal-statement}

Let micro-process be \(X_t\) and candidate macro-process be
\(Z_t = g(X_t)\). For horizon \(L\), closure asks whether \(Z_t\) is
sufficient for forecasting \(Z_{t+1}^{L}\) under a fixed regime and
admissible intervention class.

The horizon object here is the \(L\)-step path distribution, not only
one-step prediction. One-step tests can be used as diagnostics, but
objecthood claims are indexed to the declared horizon target.

Informationally, the predictive side asks whether adding \(X_t\) beyond
\(Z_t\) yields material gain for macro-future prediction.
Interventionally, the causal side asks whether macro-transition laws
remain stable across admissible perturbations without needing
within-class micro-identification.

The two sides are readings of one target, transition autonomy at the
macro level, and do not compete as separate tests.

\subsection{Why This is More Than
	Compression}\label{why-this-is-more-than-compression}

Compression is necessary but not sufficient for objecthood. A compressed
code can hide transition-relevant heterogeneity and still score well on
narrow tasks. Closure blocks that loophole by asking whether hidden
heterogeneity matters for onward macro-behavior.

Closure therefore functions as an anti-gerrymandering condition.
Convenient coding alone does not satisfy it; a coarse-graining satisfies
it when it carries stable transition structure under the declared
constraints.

This paper accepts Dennett's claim that compression
success is evidence of objective patterning, but denies that compression
success is the full ontological test \citep{dennett1991}. It also accepts
Ladyman and Ross's concern that real patterns should be projectible and
structurally disciplined, while adding an explicit transition criterion
for adjudicating borderline cases \citep{ladyman2007}. The intent is to
keep those insights while making their realist force less permissive.

The same condition answers a familiar Haugeland-style worry. The paper
does not infer objecthood directly from ``there is a useful pattern in
the data.'' It adds an explicit transition-autonomy condition that
determines when pattern talk has ontological force rather than merely
descriptive convenience \citep{haugeland1998}.

\subsection{Admissibility and Hierarchical
	Evaluation}\label{admissibility-and-hierarchical-evaluation}

Closure is always relative to an intervention class, which invites the
standard objection of circularity: who picks the class? On the present
view, the system's own dynamics do. A regime
is a region of dynamical possibility marked out by physical features of
those dynamics, such as characteristic timescales, energy scales, and
the control channels actually present, and these can be identified
before any candidate partition is scored
(\S\ref{regime-individuation-and-circularity}). The admissible
intervention class is the set of perturbations that regime tolerates
without dissolving. The analyst's role is to read off that class honestly, not to invent it. Procedural discipline (below) makes the reading auditable,
but the constraint itself comes from the dynamics, not from the
analyst. Objecthood, on this view, is stability of macro-transition
structure across the space of perturbations the regime actually
sustains, not merely fit to observed trajectories.

Admissibility is fixed upstream of objecthood verdicts by three
constraints, which make that space explicit:

\begin{enumerate}
	\item
	      Epistemic admissibility: interventions are measurable and
	      implementable by bounded agents.
	\item
	      Dynamical admissibility: interventions preserve the target regime
	      class.
	\item
	      Explanatory admissibility: interventions target variables with
	      cross-context generalizability and non-trivial counterfactual reach,
	      rather than one-off manipulations tuned to a single episode.
\end{enumerate}

These constraints are physically anchored in available control channels
and implementation structure, not in analyst preferences about
favored ontologies. The explanatory constraint is set before closure
verdicts and does not assume in advance which candidate partition will
prove stable. This gives the interventionist notion of a ``possible
experiment'' concrete bounds; if it were left at mere conceptual
possibility, nearly any hypothetical manipulation might count
\citep{woodward2015}. It also aligns with \citet{woodward2021}'s argument
that correct variable choice requires satisfying
``proportionality'': minimizing the omission of dependency relations.
Domain expertise still shapes regime identification and
control-channel selection, but that shaping is declared, auditable, and
separated from partition scoring.

Operationally, explanatory admissibility can be audited without looking
at which partition wins. A candidate intervention class is stronger when
it is implemented through control channels that already exist in the
domain and are replicable across sites, operators, or runs. It is weaker
when success depends on narrow, one-off tuning that fails under small,
feasible perturbations of the same intervention family.

Pearl defends the $do(x)$ operator
as an ideal ``surgery'' that can be defined mathematically regardless of
physical feasibility \citep{pearl2000}. For causal inference, that
idealization is productive: it lets one reason about effects even when
direct manipulation is unavailable. But macro-objecthood, on the view
defended here, requires tighter constraints than logical derivation. A partition that achieves
closure only under physically unrealizable interventions, ones that no
bounded agent could implement without destroying the target regime, does
not support the kind of stable, intervention-guiding structure that
warrants realist commitment. The admissibility constraints above ensure
that closure is assessed relative to control channels that are actually
available in the regime, not relative to arbitrary mathematical
surgeries. The restriction leaves Pearl's framework intact for causal
inference and limits only which of its outputs carry ontological weight.

We often speak of ``interventions'' as laboratory actions performed by
human scientists, but the admissible class is ultimately defined by the
objective physical ecology of the system. A system's boundaries are
constantly probed by naturally occurring environmental perturbations,
such as thermal noise, chemical diffusion gradients, and mechanical
impacts, so in this sense nature intervenes on itself. A
cell membrane objectively screens out certain molecular fluctuations
whether a human experimenter is observing it or not. The transition
dynamics of the cell carve out a regime, and that regime determines
which perturbations are admissible: those the system processes or
resists while maintaining its macro-transition structure. The
intervention class is therefore not a human invention; it is a formalization of the space of
possibility the dynamics already define.

Evaluation proceeds in a fixed order.

\begin{enumerate}
	\item
	      Fix regime, horizon, admissible intervention class, diagnostics, and
	      model class.
	\item
	      Compare candidate partitions under those fixed constraints.
	\item
	      If constraints are revised after inspecting results, restart and
	      disclose the revision.
\end{enumerate}

For admissibility, this means that intervention classes are fixed first,
from independent physical constraints and established control practice
in the target domain, and closure is then evaluated only relative to
that fixed class. The order blocks back-fitting and makes disagreement
tractable. If verdicts remain highly sensitive across nearby defensible
admissibility specifications, commitment should be downgraded.

\subsection{Regime Individuation and
	Circularity}\label{regime-individuation-and-circularity}

If regime individuation already commits to which macro-variables
matter, the protocol risks a version of circularity: closure testing
presupposes regime identification, but regime identification might
presuppose partition outcomes. In reply, defensible regime
specifications draw on quantities characterizable without presupposing
which candidate partition will succeed. Temporal parameters, energy
scales, flow rates, and enforcement intensities are typically
identifiable through physical indicators that competing candidate
partitions can agree on before either is evaluated. In the traffic
illustration, ``weekday commuter traffic in one metropolitan corridor''
is characterized by clock schedules, total vehicle counts, and road
network topology, all quantities that investigators favoring either the
lane-flow or vehicle-ID-cluster partition can agree on before either
closure test is run. Regime specifications of this kind are anchored in
lower-level or cross-cutting physical indicators rather than in the
outcome of the test they enable. Regime specifications should be
grain-neutral across competing candidate partitions. A regime can itself be a macro-level description while
still not presupposing which finer-grained partition will close
within it.

The independence is imperfect in some domains, particularly biology and
social science, where macro-variable choice and regime characterization
can be entangled. The present account treats that entanglement as a
source of evidential weakness rather than a concealed defect. When
alternative, independently motivated regime specifications produce
divergent closure verdicts, the verdict ceiling should be downgraded to
at most qualified. Regime individuation and partition evaluation are
iterative rather than viciously circular. Initial regime characterizations are provisional, and closure
testing can refine them. The process stabilizes into reflective
equilibrium in the sense Goodman articulated: rules of inference and
particular inferential judgments are justified by mutual coherence
rather than by either serving as the unilateral foundation for the
other \citep{goodman1955}. If iteration does not stabilize across nearby
defensible starting specifications, the result should be treated as
indeterminate for that regime rather than forced into a robust verdict.
Here stabilization means that comparative ordering across independently
specified candidate partitions stops shifting under small, defensible
protocol variations; it does not presuppose one privileged grain in
advance.
The discipline comes from requiring that all revisions be disclosed and
that each revision restart the evaluation rather than silently inheriting
prior results.

A more global form of the worry asks what selects the right regimes and
contexts in the first place. Some sufficiently narrow regime can
probably be found in which almost any candidate looks stable, so it
might seem that closure needs a privileged scientific task to say which
contexts count. It does not, because the protocol checks two kinds of
robustness that do different jobs.

Robustness across the admissible intervention class is part of what
closure is. The criterion asks whether within-class micro-differences
alter macro-transitions under the perturbations in that class, so a
candidate that holds up under some admissible perturbations and fails
under others is not closed under that class. That failure is a fact
about the candidate, not a gap in the evidence.

Robustness across nearby regimes, horizons, and modelling choices plays
a different role. It is not part of closure in a single regime, since a
candidate can be closed in one regime and leak in the next, as ice does
across a phase boundary. It is evidence that the regime under test is a
real feature of the system's dynamics rather than one drawn to suit
a favored partition. A regime cut so narrowly that a favored candidate
looks stable within it, and nowhere nearby, will usually fail
projectibility under small shifts in context
(\S\ref{regime-dependence-and-projectibility}), and a regime whose
boundaries were drawn by looking at closure scores violates the
evaluation order of \S\ref{admissibility-and-hierarchical-evaluation}.
A closure verdict is therefore always relative to a regime, but the
regime has to be identified from the system's physical features and has
to hold up under nearby variation. No further task-level justification
of the context is required, because the contexts that matter are the
ones in which the system actually supports autonomous transition
structure.

\subsection{Pattern Reality Versus
	Macro-Objecthood}\label{pattern-reality-versus-macro-objecthood}

A pattern can be real in a weaker sense without satisfying this
paper's objecthood criterion. A representation can capture regularities
that are descriptively useful while still failing closure under
admissible interventions. The criterion here is stricter because it
targets macro-objecthood, not every projectible summary. The paper can therefore
preserve a generous attitude toward pattern detection while denying
ontological standing to high-leak candidates.

The distinction also bears on composition debates. In van Inwagen's terms, the standing question is
when some things compose a further thing \citep{vaninwagen1990}. The
present proposal does not offer a full mereology, and it recasts the
question rather than answering it on its own terms. What makes a
macro-object an object, on this view, is that a pattern carries its own
transition structure at its grain. The lower-level constituents realize
that pattern; they are not what its objecthood consists in. Composition
claims are then disciplined by the same condition: a claim that some
parts make up a macro-object is warranted when the pattern they realize
is robustly closed under a declared regime, horizon, and intervention
class, and should be withheld
when candidate sums remain transition-fragile. This leaves the view
compatible with several positions on composition, including
structuralist views on which talk of objects can be given up in favor
of talk of structure \citep{french2014}. One disagreement does remain.
The criterion treats closed macro-level transition structure as a real
feature of the world. A view that denies this, and holds that only the
micro-level facts are real, disagrees with the paper about substance. A
view that grants it but declines to call the resulting structures
``objects'' disagrees only about the word.

This fixes the dialectical burden. The paper need not show that
non-closed representations are worthless, only that they do not meet
the objecthood standard defended here.

\subsection{Admissibility Disputes}\label{admissibility-disputes}

Admissibility disputes are expected, especially across domains. The
adjudication rule is to compare candidate intervention classes by what
they physically permit and whether they preserve the same target regime.
Boundary-pressure perturbations can be admissible for storm dynamics,
while molecule-by-molecule remote rewriting is not. In institutional
settings, changing enforcement intensity can be admissible, while
instant arbitrary rewriting of all agent commitments is not. The choice
of class bears directly on the verdict. A storm satisfies closure under
boundary perturbations because its pressure organization and rotational
structure determine macro-transitions without tracking individual
molecules. Under molecule-by-molecule rewriting, the same candidate
fails, because admissibility now permits interventions that reach inside
the macroclass and exploit within-class differences. The candidate has
not changed; the admissible control channel has, and the verdict changes
with it. Admissibility must therefore be declared and justified before
scoring, and not adjusted afterward to protect a preferred outcome.

For transparency, three template classes are often useful as a starting
grid: boundary nudges, coarse actuator controls, and policy levers.
These templates are not universal, but they make admissibility
comparison public and repeatable.

Interventions that directly target preservation of partition labels,
rather than system variables, are inadmissible without exception. They
would secure closure by stipulation, whereas the test asks whether
transition structure is autonomous under physically meaningful controls.

When two intervention classes are both admissible and target the same
regime, the criterion permits plural testing rather than forced monism.
Verdicts should report which admissibility class was used and how
sensitive results are across nearby defensible classes. A candidate partition that appears closed only under one narrowly tuned
class should be downgraded when its results shift across nearby
defensible alternatives.

\subsection{Canonical Criterion
	Statement}\label{canonical-criterion-statement}

The core criterion can be stated compactly.

A candidate coarse-graining \(Z=g(X)\) qualifies for macro-objecthood in
regime when, under fixed horizon \(L\) and admissible intervention class
\(\mathcal{I}\), macro-transition structure is autonomous up to declared
tolerance, so within-class micro-differences do not materially improve
macro-future prediction or intervention-guided control.

Candidate partitions must be specified independently of the particular
sample trajectory used to score closure, for instance by a pre-declared
construction rule or learning objective fixed before evaluation.
Otherwise closure collapses into bespoke encoding, where any dataset can
be made to look autonomous by post hoc recoding.

In exact Markov settings, strong lumpability is a sufficient benchmark
realization of this criterion. In non-ideal settings, leakiness-centered
diagnostics estimate distance from that ideal under fixed constraints.

\subsection{Predictive and Interventional
	Closure}\label{predictive-and-interventional-closure}

Predictive and interventional closure should be distinguished but not
separated. Predictive closure concerns whether macrostate information
screens off transition-relevant micro-detail for macro-future
forecasting. Interventional closure concerns whether macro-transition
structure remains stable under admissible perturbation.

The evidential relation is asymmetric. Strong interventional closure
typically implies predictive adequacy for the same target and horizon,
while predictive adequacy alone can persist in cases where
interventional stability fails under distribution shift. For that reason,
the criterion treats predictive evidence as important but not final.

The asymmetry also sets the burden of proof. Claims to robust
macro-objecthood need either direct interventional support or compelling
indirect evidence that tracks intervention-relevant invariance. Without
either, the responsible verdict is qualified or indeterminate.

\subsection{Why Closure Carries Ontological
	Weight}\label{why-closure-carries-ontological-weight}

A reviewer can still ask why transition autonomy should count as
objecthood, and not merely as a success condition for modeling. The
reply rests on the paper's explicit commitments. Under structural
realism, ontology should track stable relational and causal
organization, not intrinsic micro-identity as such. Under
interventionism, causally relevant structure is structure that supports
stable manipulation and control.

The meta-philosophical stance here follows what Woodward calls the
methodological path to ontology: causal-interventionist tools are used
to determine what earns realist commitment, rather than seeking
independent metaphysical ``truth-makers'' or ``grounds'' that stand
apart from the practices of manipulation and prediction \citep{woodward2015}. This is not a retreat to instrumentalism; the claim is that
scientific ontology, disciplined by closure, is the only ontology
required for macro-objecthood verdicts. The demand for a further
metaphysical foundation beyond stable dynamical autonomy under
admissible intervention mistakes a philosophical habit for a substantive
requirement.

A persistent metaphysical worry is that prediction and interventional
control are merely epistemic lenses, leaving the underlying ontology
untouched. This confuses the test with what is being tested. Prediction
and intervention are diagnostic tools that reveal pre-existing
statistical structure. The conditional independencies that allow a
macro-state to screen off micro-details exist whether or not anyone
computes them. Locality, finite signal propagation, and thermodynamic
limits create these independencies as physical facts. Our predictive
models succeed when they align with this objective causal topology, and
fail when they do not. Closure is therefore an ontological criterion
because it identifies where the world's phase space objectively
decouples, and not merely where our minds prefer to chunk it.

Given those commitments, closure is the criterion that identifies when a
candidate macro-description tracks stable structure at the level where
explanation and intervention are being assessed. It is not a convenient
heuristic layered on top of ontology. On this basis the view can remain
compatible with microphysical completeness while still defending
non-redundant macro commitment.

\subsection{Structural Contrasts}\label{structural-contrasts}

Pearl argues that causal models work because nature consists of
autonomous mechanisms: changing one structural equation need not disrupt
the others \citep{pearl2000}. Closure is related to this property but not identical to
it. Pearl treats modularity as a structural assumption about a given
causal model; the present criterion asks a prior question, namely
whether a candidate coarse-graining has earned the right to appear as a
variable in such a model at all. A partition that requires persistent
micro-bookkeeping to preserve its transition structure has not isolated
an autonomous mechanism at the macro grain in the relevant sense. So
closure can be understood as an ontological precondition for the kind of
modularity Pearl's framework presupposes, rather than as a restatement
of it.

Because the criterion is modal and dynamical, and not merely
extensional, it is not vulnerable to standard ``bare structure''
worries. It concerns counterfactual transition organization under
interventions, not abstract isomorphism alone.

Mechanistic accounts of levels are right that organized entities and
activities can support stable explanation and intervention \citep{machamer2000,craver2007}. But by themselves they do not settle whether a proposed
macro-variable has been carved at the right place. Closure addresses
that prior sorting question. A candidate may be
heuristically useful or mechanistically suggestive and still fail
objecthood if admissible perturbation reveals continued dependence on
transition-relevant microstructure.

This also addresses Newman-style triviality pressure in structural
realism, where purely extensional structure can be cheap under
redescription \citep[pp.~141--145]{ainsworth2009}.
Closure requires transition and intervention stability under declared
constraints, not extensional fit alone, and arbitrary recodings
typically fail that requirement.
When apparent success depends on privileged representational accidents,
that encoding dependence is itself evidence of leakiness.

It might seem enough to call a closure-supporting partition a good
variable. But a good variable can be merely convenient, and the worry is
then that closure only tracks inferential convenience rather than what
is objectively there. The present account rejects that deflation
because closure claims are indexed to intervention classes, not only to
passive prediction. A representation
can score well on retrospective fit by exploiting accidental
correlations and still fail when admissible interventions shift
transition pathways. A closure-supporting partition survives those
shifts because the screened-off structure is not accidental. If a
candidate supports reliable counterfactual control and regime-stable
transition laws, it is tracking objective modal structure: stable
counterfactual dependencies and lawlike transition regularities under
admissible intervention. On the paper's commitments, treating that
structure as merely representational would undercut the very realism the
criterion is designed to preserve. The ontological burden is carried by
the modal profile and not by in-sample fit. Predictive sufficiency can
reveal compression, whereas interventional stability identifies
difference-makers, because it tests whether the same macro-transition
structure survives admissible manipulation instead of merely fitting
observed trajectories.

Finite agents will often have partial, noisy, and regime-limited
evidence. That calls for epistemic humility, not ontological deflation.
The uncertainty does not collapse the difference between dynamically
autonomous and dynamically incoherent partitions. That difference can be
difficult to estimate, but it is still a difference in the world rather
than in preference.

\subsection{Levels of Claim and
	Representation}\label{levels-of-claim-and-representation}

Some recurring misunderstandings come from sliding between different
levels of claim. The paper keeps four levels apart throughout.

\begin{enumerate}
	\item
	      World dynamics: the implemented process itself.
	\item
	      Pattern type: stable transition structure supported by that process.
	\item
	      Pattern token: a concrete instance under specific boundary conditions.
	\item
	      Representation: a model, equation, classifier, or narrative that
	      tracks the pattern.
\end{enumerate}

A representation can fail while the pattern type remains real. A token
can fail while the type remains robust. A token can also succeed while
the type is weak, for example in narrow calibration windows. Closure
claims in this paper are primarily type-level claims under specified
regimes, then secondarily claims about token reliability under
perturbation. The criterion therefore does not equate model performance
with ontology. It asks whether the represented transition structure is
autonomous, not whether one representation currently performs well.

\subsection{Closure, Underdetermination, and Objecthood
	Discipline}\label{closure-underdetermination-and-objecthood-discipline}

A remaining concern is underdetermination. Even after the criterion is
stated, a reviewer may argue that many distinct coarse-grainings can be
tuned to look acceptable on available data. If so, closure might seem to
collapse back into model choice rather than objecthood discipline.

Three questions that are often conflated need to be kept apart.

\begin{enumerate}
	\item
	      Which candidate partitions are descriptively serviceable in a dataset?
	\item
	      Which candidates remain transition-autonomous under admissible
	      perturbation?
	\item
	      Which of those candidates remain stable under modest horizon and
	      admissibility variation?
\end{enumerate}

Underdetermination is strongest at the first question and weaker at the
second. It is often weakest at the third. Many candidates can fit
observed trajectories. Far fewer preserve counterfactual structure when
the regime is probed. Fewer still remain stable across nearby defensible
setups. This staged filtering, which goes beyond fit alone, is where
closure earns its metaphysical role.

The account does not claim that every domain yields one uniquely
privileged partition. Its claim is that objecthood commitments should be
constrained by transition autonomy and robustness, rather than by fit
alone. When underdetermination persists after those filters, the
responsible result is qualified or plural commitment under transparent
constraints, not forced monism and not unconstrained relativism. When
exact closure is unavailable, leakiness-centered diagnostics are used to
estimate comparative distance from closure, and commitment strength
tracks the stability of that evidence.

\section{Exact Benchmark: Strong
  Lumpability}\label{exact-benchmark-strong-lumpability}

Strong lumpability is introduced as a benchmark, not as a condition
that every macro-object must meet exactly. It gives a clean exact case in Markov
settings where closure can be stated and checked without ambiguity.

Let partition cells under \(g\) be macroclasses. Strong lumpability
holds when microstates within the same macroclass induce identical
transition probabilities to all macroclasses \citep{kemeny1960}. When this condition holds,
macro-transitions are autonomous by construction.

More precisely: for any two microstates \(x, x'\) in macroclass \(C_i\),
and any macroclass \(C_j\), strong lumpability requires
\(\sum_{y \in C_j} P(y \mid x) = \sum_{y \in C_j} P(y \mid x')\).

\textbf{Lemma.} If the micro-process is first-order Markov and the
partition induced by \(g\) is strongly lumpable, then the induced
macro-process is itself Markov and transition-autonomous at the macro
level.

\emph{Proof sketch.} Define a macro-kernel by summing micro-transition
probabilities from any representative \(x \in C_i\) into each macroclass
\(C_j\). Strong lumpability guarantees this sum is
representative-independent. The macro-kernel is therefore well-defined,
and the induced process over macroclasses is Markov with transitions
given by that kernel.

A partition that satisfies this condition preserves transition
structure at the macro grain, so its standing does not rest on
usefulness alone. A partition that fails it mixes
transition-heterogeneous microstates and therefore lacks macro autonomy.

Computational mechanics gives an equivalent formulation, which shows
that the condition is more than superficial bookkeeping. For any
macro-process \(Z\), two prediction
machines can be defined. The \(\varepsilon\)-machine is the minimal
model that predicts \(Z\)'s future from \(Z\)'s own past, using only
macro-level information. The \(\upsilon\)-machine is the minimal model
that predicts \(Z\)'s future from the full micro-past \(X\), using
everything available \citep{shalizi2001,rosas2024}.
Closure holds when these two machines are equivalent: the
\(\varepsilon\)-machine already captures everything the
\(\upsilon\)-machine knows about macro-futures, and extra
micro-information is redundant for the macro target. In this exact
ideal, closure holds without approximation.

Exact strong lumpability is uncommon in open, noisy, and
path-dependent systems. That limitation motivates the graded extension;
it does not undermine the criterion.

\subsection{Minimal Formal
	Illustration}\label{minimal-formal-illustration}

The anti-gerrymandering role can be shown with a minimal symbolic
example. Let microstates be \(\{x_1,x_2,x_3,x_4\}\) with transition
rows:

\[
	\begin{aligned}
		P(x_1,\cdot) & = (\alpha,\beta,0,0),  \\
		P(x_2,\cdot) & = (\alpha,\beta,0,0),  \\
		P(x_3,\cdot) & = (0,0,\gamma,\delta), \\
		P(x_4,\cdot) & = (0,0,\gamma,\delta),
	\end{aligned}
\]

with \(\alpha+\beta=1\) and \(\gamma+\delta=1\).

Partition \(A\) groups \(\{x_1,x_2\}\) and \(\{x_3,x_4\}\). Partition
\(B\) groups \(\{x_1,x_3\}\) and \(\{x_2,x_4\}\). In \(A\), members of
each macroclass have matching onward profiles to macroclasses, so
macro-transitions are autonomous. In \(B\), members of one macroclass
differ in onward profiles whenever \(\alpha \neq \gamma\), so the
macro-label hides a transition-relevant difference.

Partition \(B\) can still be made predictive by adding repeated
within-class bookkeeping. A partition is \emph{high-maintenance} when
reliable macro-prediction repeatedly forces refinement of within-class
distinctions, so that yesterday's macrostate labels must be supplemented
by new micro-bookkeeping to retain accuracy. This paper excludes such
partitions from robust macro-objecthood.
On a fixed sample trajectory, this bookkeeping can make Partition \(B\)
look nearly as predictive as Partition \(A\). The difference appears
under admissible perturbation: \(A\) preserves macro-transition structure
without renewed microstate tracking, whereas \(B\) does not.

Both partitions are definable, but only the first preserves transition
autonomy. Closure therefore discriminates between legitimate
macro-candidates and merely codable aggregates.

\subsection{If an Unusual Partition
	Closes}\label{if-an-unusual-partition-closes}

A common reaction is that a strange disjunctive partition might satisfy
the formal condition in a symmetric system. It can. On this account, if
such a partition supports autonomous macro-transitions under fixed
constraints, it counts as real at that grain.

Accepting this concedes nothing to arbitrariness, because the criterion
tracks transition coherence and not intuitive naturalness. If one wants
a separate naturalness filter, that is an additional criterion and
should be declared as such.

The exclusion claim should therefore be read narrowly. The criterion
excludes dynamically incoherent, high-leak aggregates, not every unusual
grouping.

\subsection{What the Formal Machinery Alone Does Not
	Provide}\label{what-the-formal-machinery-alone-does-not-provide}

Formal closure is a mathematical property; the criterion defended here is a philosophical one. Four elements of the present proposal have
no counterpart in the formal literature alone.

First, intervention indexing. Causal-state constructions classify
histories by equivalence of future distributions under the observed
process, which is primarily an observational-predictive equivalence
relation \citep{shalizi2001}. The present criterion adds an
admissible intervention class as a parameter of the closure test.
Objecthood claims are indexed to what the system does under
perturbation as well as under passive observation, and this separates
the criterion from a purely statistical diagnostic.
Rosas et al.'s framework distinguishes informational, causal, and
computational closure as complementary formal diagnostics
\citep{rosas2024}. The present paper treats predictive sufficiency as
evidential, but ties objecthood commitment to admissible-intervention
stability and to the absence of hidden micro-identity bookkeeping.

Second, an explicit ontological interpretation. Strong lumpability tells
you when a coarse-graining preserves transition structure. It does not
tell you whether that preservation warrants realist commitment. The
philosophical work is to argue that, under stated commitments,
transition autonomy under admissible interventions is sufficient for
macro-objecthood, and to give the conditions under which that claim
should be downgraded or withdrawn.

Third, a graded commitment protocol. The formal literature offers exact
conditions and, more recently, approximate-closure measures. It does not
supply a protocol for translating those measures into warranted
ontological verdicts with explicit robustness requirements and downgrade
rules. The selective commitment structure (robust, qualified,
indeterminate) is a philosophical addition.

Fourth, an anti-gerrymandering argument. The claim that closure is
the right anti-gerrymandering condition, and that compression without
closure is insufficient for objecthood because compression can be won by
codings that effectively preserve hidden micro-identities, is a
philosophical claim rather than a theorem. Sections 2 through 4 defend it.

Intervention indexing matters most where observational and
interventional verdicts come apart. An aggregate
macro-variable can track a latent process so faithfully that it achieves
observational predictive sufficiency: an \(\varepsilon\)-machine built
on the proxy compresses past observations and predicts future ones as
efficiently as the full-state representation. Classical computational
mechanics would recognize it as identifying a real pattern. But under
admissible intervention on the proxy itself, transition structure
collapses, because the variable was tracking a common cause rather than
instantiating an autonomous mechanism.

Suppose, for instance, that a composite variable aggregates
two subsystems that are both driven by a shared upstream oscillator. Under
passive observation, the composite's macro-state predicts its own future
as well as the full micro-state does, because the shared driver
synchronizes the subsystems and makes their joint trajectory
compressible. An \(\varepsilon\)-machine built on the composite would
find no residual micro-information worth adding. Yet if an admissible
intervention targets the composite directly, say by perturbing one
subsystem while leaving the driver intact, the previously synchronized
trajectories decouple and the composite's macro-transition structure
breaks down. The predictive sufficiency was parasitic on the common
cause; the composite itself had no autonomous transition organization.
Observational diagnostics alone cannot distinguish this case from genuine interventional closure. The present criterion can, because it indexes objecthood to
intervention-stability and not to observational compression alone.

The case invites a worry from the other direction. Organisms,
artifacts, and institutions all depend on their surroundings: on
gravity, on a temperature range, on nutrient flows, on the hand that
swings the hammer, on supporting practices. If any dependence on
structure outside the partition counted against objecthood, the
common-cause case would condemn nearly every special-science kind. The
line is drawn by distinguishing two ways an environment can support a
candidate.

Some environmental conditions are held fixed by the regime: gravity,
the temperature range, the institutional background. They are part of
the regime's specification. The candidate's closure is assessed against
them, and a change in them is a change of regime rather than a
perturbation within it. Other influences vary and act on the candidate
as inputs, as the swing acts on the hammer or a nutrient supply on a
cell. These are treated as exogenous inputs that the transition law
takes as given, not as part of the candidate's state. The closure
question is then whether the candidate's own macrostate, together with
those inputs, fixes its macro-future. For the hammer and the cell it
plausibly does. Given the swing, the hammer's position and momentum
settle what happens to the nail, and the lump's microstructure adds
nothing. Dependence of this kind leaves the candidate's own state doing
the predictive work.

The composite driven by the shared oscillator is different. Once the
driver's state is given, the composite's own present state adds nothing
to the forecast of its future, since the driver sets both subsystems.
Its present state looks predictive only because, in the observed
process, it happens to mirror the driver. Its apparent macro-dynamics
are the driver's dynamics seen through a correlate, and an admissible
intervention on the composite breaks the correspondence while leaving
the regime intact. Since everything depends on its environment, the
question for any candidate is whether its own macrostate carries the
forecast or only stands in for an off-partition variable that the
observed process keeps aligned with it. An organism
kept alive by a stable environment does not fail closure for that
reason alone. It would fail only if an admissible perturbation acting
on it showed that its apparent macro-dynamics had tracked an external
driver all along.

A similar point applies to causal-emergence work. Effective-information
analyses identify when macro-descriptions gain determinism relative to
micro-descriptions \citep{hoel2013}; \citet{dewhurst2021} argues that
such gains support only epistemic emergence absent further constraints.
That diagnosis is correct as far as it goes, but it identifies a gap
without filling it. A partition can increase effective information while remaining
fragile under intervention or horizon variation. The present approach
treats such gains as corroborating evidence within a closure-governed
criterion, not as a standalone ontology test.

The resulting position borrows formal tools from computational mechanics
and causal-emergence analysis, but it is not a relabeling of either. It
adds intervention-indexed objecthood conditions, an explicit commitment
protocol, and a sustained argument about why closure is the right fix
for pattern-realism's permissiveness problem.

\section{Approximate Closure Without Instrumentalist
  Drift}\label{approximate-closure-without-instrumentalist-drift}

\subsection{Why Approximation is
	Expected}\label{why-approximation-is-expected}

In complex systems, boundaries leak and couplings shift across regimes.
Exact closure is therefore unusual. A credible realism criterion cannot
require perfect closure everywhere. The study of approximate
aggregation has a long history, especially Simon's account of hierarchy
and near-decomposability, where strong within-subsystem couplings can
coexist with weaker cross-subsystem couplings and thereby separate
short-run local dynamics from longer-run aggregate behavior
\citep{simon1962,ando1961}. The present extension is philosophical: it
uses graded closure to support graded ontological verdicts, and it does
not aim to improve estimation. Provided constraints are fixed and
diagnostics are comparative, approximate closure is the expected
non-ideal form of the same criterion and concedes nothing to
arbitrariness.

The right ontology test is type-level before token-level. A pattern type
can be robustly closed for a regime even when a specific token fails
because of boundary violation, atypical perturbation, or timescale
mismatch. One anomalous token is therefore not decisive. The relevant
question is whether failures are exceptional or systematic for the type
under the stated constraints. Systematic failure, where within-class
micro-differences repeatedly matter for macro-transitions across tokens
under fixed constraints, triggers type-level downgrade rather than
token-level exception.

\subsection{Leakiness as Canonical
	Target}\label{leakiness-as-canonical-target}

Leakiness measures how much within-class micro-detail still improves
macro-future prediction once current macrostate is fixed. A canonical
quantity is the conditional mutual information \(I(X_t; Z_{t+1} \mid Z_t)\)
\citep[eqs.~(2.60)--(2.61)]{cover2006}. Low values support closure.
High values indicate hidden transition-relevant heterogeneity.

Using information-theoretic measures to define ontology might appear to
assume a computationalist metaphysics, but the justification is
strictly physical. Information processing incurs thermodynamic costs
\citep{landauer1961}. A system, whether a biological organism or an
engineered control mechanism, that tracks full micro-details without
gaining macro-predictive advantage incurs unsustainable energetic
overhead. Evolution and physical self-organization therefore robustly
select against leaky boundaries. Measuring leakiness is thus a formal
accounting of the physical and thermodynamic constraints that force
systems to dimensionally reduce their own state spaces, and not a
computer-science metaphor imposed on nature. A low-leakage partition
tracks the joints where the physical world successfully minimizes this
metabolic and computational cost.

In this paper, leakiness is the default target quantity. Other
diagnostics, such as within-class transition divergence and predictive
gain from added micro-features, function as estimators or proxies when
direct estimation is limited.

Leakiness should be read comparatively and through robustness
checks, not as a context-free absolute cutoff. Absolute tolerance is
domain-relative, but high sensitivity of leakiness rankings to modest
defensible modeling choices is itself a downgrade trigger.

Comparative ranking is necessary but not sufficient. A candidate can
beat nearby alternatives and still fail objecthood if absolute leakiness
remains high and instability persists under modest robustness checks.
The criterion therefore requires both relative superiority and minimum
viability.

The viability floor can be stated without a numeric threshold. A
partition fails minimum viability when, across the declared robustness
checks, within-class micro-features yield consistent, non-negligible
gains in macro-future prediction or intervention response relative to
the best macro-only model. Persistence of that pattern across
diagnostics and perturbation tests is decisive. The partition has not
achieved the transition autonomy the criterion requires, regardless of
how it ranks against competitors.

The term ``non-negligible'' can be given a comparative anchor without
fixing a universal numeric threshold. Leakiness is most clearly
non-negligible when adding within-class micro-features yields prediction
gains comparable in magnitude to the gains the macro-partition itself
provides over a naive baseline model. If the macro-partition delivers
substantial predictive improvement over no partition at all, and
within-class micro-features add only marginal further gain, that
residual leakiness is negligible for most objecthood purposes. If
micro-features instead add gains approaching or matching the
macro-partition's own contribution, the partition is not doing most of
the relevant work, and minimum viability is not met. This framing turns
the judgment into one about a single tractable quantity, relative
predictive contribution, instead of an abstract absolute cutoff, and it
makes cross-domain comparisons more tractable. The two requirements work
in sequence. The viability floor is applied first as a minimum gate,
excluding candidates regardless of how they compare to rivals, and
comparative ranking then assigns verdict categories among candidates that
pass it. To keep this anchor non-discretionary, the baseline model class
and gain metric must be fixed upstream with regime and admissibility,
and post hoc revisions trigger restart and disclosure.

Verdict-category boundaries are vague because the underlying property
of transition autonomy is itself graded. Biological fitness is similar:
it is real and explanatorily significant without admitting a sharp
boundary between fit and unfit. What matters primarily is ordering
stability. When partition
$A$ consistently shows lower leakiness and greater perturbation
stability than partition $B$ across diagnostics, that ordering should be
robust even if the exact placement of the boundary between ``qualified''
and ``robust'' carries some domain-conventional element. A partition
that dominates alternatives across diagnostics and remains stable under
modest changes to horizon, intervention distribution, and robustness
checks earns the stronger verdict wherever one draws the
categorical line.

\subsection{Procedural Safeguards in Non-Ideal
	Cases}\label{procedural-safeguards-in-non-ideal-cases}

Approximate closure claims require explicit safeguards.

\begin{enumerate}
	\item
	      Fix regime, horizon, admissibility class, and diagnostics before
	      comparative scoring.
	\item
	      Compare candidate partitions under the same fixed setup.
	\item
	      Test robustness under modest changes in horizon and intervention
	      distribution.
	\item
	      Apply a minimum-viability floor: if all candidates remain high-leak
	      and fragile, return no robust objecthood verdict for that regime.
	\item
	      Report verdicts as robust, qualified, or indeterminate.
\end{enumerate}

These safeguards also fit Wimsatt's robustness-analysis idea:
confidence should rise when partially independent ways of detecting the
same structure converge, and it should fall when apparent agreement is
driven by shared assumptions or disappears under small perturbation
\citep{wimsatt1981}. Multiple diagnostics do not replace criterion-level
argument with a vote count. They reduce the risk that a favorable verdict
is an artifact of one estimator, one horizon choice, or one admissibility
specification.

The safeguards keep threshold choice procedural rather than
discretionary. They also locate estimation difficulty correctly, as an
epistemic limit on access and not a defect in the metaphysical criterion
itself.

\textbf{Selection discipline.} A natural concern is that flexibility
can silently re-enter through upstream choices. One rule addresses it:
admissibility, state construction, horizon, and diagnostics are fixed
before scoring, and post hoc revisions trigger restart and disclosure.
The rule does not eliminate judgment, but it makes judgment auditable
and prevents favored partitions from being protected by moving criteria.

\subsection{Regime Dependence and
	Projectibility}\label{regime-dependence-and-projectibility}

Regime dependence does not imply observer-relativity. Closure depends on
actual transition structure under specified constraints, and a pattern
can be closed in one regime and leaky in another because the structure
has changed, or because the available control
channels have shifted. Since the admissible intervention class is defined by the
regime's transition structure
(\S\ref{admissibility-and-hierarchical-evaluation}), a change in that
structure can alter both the admissible class and the closure verdict even when the
system's micro-dynamics remain the same. Regime and horizon index the
modal profile being claimed; they are not a concession that anything
goes.

Projectibility provides a further check. Defensible regime
specifications should support stable induction under modest
counterfactual variation. Narrowly engineered regimes that protect a
favored partition usually fail under small context shifts. When that
occurs, the candidate was never robustly closed in the relevant sense.

Here the paper is closest to Ladyman and Ross. Projectibility functions
as a realism-relevant stress test on whether a closure claim survives
beyond calibration conditions, rather than as an optional methodological
virtue \citep{ladyman2007}. A closure claim that lacks projectible
robustness cannot support a robust objecthood verdict.

Regime dependence itself takes two forms. Ontic regime
dependence occurs when system structure changes, such as phase
transitions or boundary-condition shifts. Epistemic regime dependence
occurs when measurement or control access changes while underlying
structure remains fixed. The first changes what is there to be tracked.
The second changes what can be warranted from available evidence.

\subsection{One Distributed
	Illustration}\label{one-distributed-illustration}

Macro does not mean large or spatially contiguous. Coarse-graining can
be logical and distributed. Monetary systems illustrate this point
without requiring new machinery.

The macro-transition structure of transactions can remain stable across
heterogeneous micro-realizations, such as cash tokens, ledger records,
and digital balances, under admissible legal and financial
interventions. If that stability holds, macro-objecthood is warranted in
regime. If implementation channels degrade, closure can fail even while
symbolic representations persist.

In a minimal failure case, a state retains formal legal tender rules
while losing reliable payment enforcement and settlement implementation.
The rules have not changed, but the physical infrastructure that
enforced screening-off is gone, so closure fails. The representation
persists but transition autonomy degrades, and closure-based commitment
should be downgraded. The same structure applies more broadly: the
criterion explains when and why a social entity stops being a real
macro-object without requiring the claim that it was never real, or that
it remains real simply because people believe in it.

Organisms give a complementary illustration. In many physiological
regimes, membrane and regulatory organization screen off large amounts of
molecular variation for the target transitions, so intervention on
organ-level variables remains predictively and manipulatively effective.
When those screening structures fail, for instance under severe systemic
breakdown, the same macro-description can become leaky and require
finer-grained tracking. Organism talk is therefore not closed in every
regime, and the example shows that closure can be regime-robust without
being regime-universal.

The three verdict categories sort these illustrations differently. The
monetary example in a well-functioning system
approaches robust objecthood in the declared regime: macro-transition
structure remains stable across heterogeneous realizations and admissible
interventions. The organism example falls into the qualified range.
Within established physiological regimes, membrane organization and
regulatory control screen off large amounts of molecular variation,
making organ-level interventions reliably effective, but leakage appears
at ecological or pathological boundaries where those screening structures
degrade. The failed-implementation case
represents a downgrade: formal representations persist while transition
autonomy erodes, warranting withdrawal of robust commitment. In each case
the verdict is an outcome of the same diagnostic criteria applied to a
different evidential situation.

These illustrations are meant only to show that closure tracks
transition structure rather than spatial shape. Nothing in the argument depends
on taking a stance on broader social ontology.

\subsection{Failure Conditions and Downgrade
	Rules}\label{failure-conditions-and-downgrade-rules}

The account should also say clearly when commitment should be
withdrawn or downgraded. Four failure patterns are especially relevant.

\begin{enumerate}
	\item
	      Persistent high leakiness across defensible diagnostics under fixed
	      constraints.
	\item
	      Strong disagreement among diagnostics that does not resolve under
	      modest robustness checks.
	\item
	      High sensitivity of verdicts to small, defensible changes in
	      admissibility or horizon.
	\item
	      Candidate sets where no partition clears the
	      minimum-viability floor.
\end{enumerate}

When these patterns appear, the verdict should be explicitly downgraded
to qualified or indeterminate status instead of being forced into a
binary judgment. The criterion does not pretend that every case must
produce a sharp ontology verdict.

Conversely, when closure persists across distinct, independently
motivated admissibility classes within the same regime, that
cross-admissibility stability raises evidential confidence for robust
commitment.

A downgraded token does not automatically refute type-level closure.
What matters is whether the instability is exceptional at token level
or systematic for the type under the stated constraints.

\subsection{Stable Versus Merely Entailed
	Patterns}\label{stable-versus-merely-entailed-patterns}

On the graded account, a pattern can be entailed by microhistory and
laws without being stable as a macro-handle. Entailment alone is cheap,
whereas stability under admissible perturbation is demanding.

The distinction bears directly on permissiveness debates, because a
contrived disjunctive construction can be true of a realized trajectory
while still failing closure. It can require continual within-class
micro repair, fail under modest regime shifts, and lose interventional
reliability.

A historical case makes the pattern concrete. Phlogiston theory
compressed combustion phenomena under a single macro-variable, but the
partition required persistent bookkeeping to survive. When metals gained
weight upon burning, theorists introduced ``negative phlogiston''; when
different substances showed different weight changes, further ad hoc
parameters appeared. Each patch was a new within-class distinction
imported to preserve macro-prediction. Oxygen theory required
substantially less such repair, since its macro-variable (oxidation state)
screened off the relevant chemistry without persistent bookkeeping. The
criticism is structural and does not rest on retrospective ridicule. A
candidate that continually imports corrections to preserve
macro-prediction behaves like a non-autonomous partition, and the
closure framework detects this failure directly.

Gerrymandered constructions need not be false. The claim is that they
usually fail to qualify as macro-objects. They may remain descriptions,
but not object-level descriptions in the relevant regime. The same
applies to computationally adequate models more broadly. A higher-level
model can be useful for prediction or control without meeting closure
standards. Computational adequacy means the model serves some purpose;
closure requires that transition-relevant micro-differences are screened
off in the specified regime. A computationally adequate but high-leak
representation remains a valuable tool, but its adequacy does not
automatically confer macro-object status.

\subsection{Practical Qualification Under Sparse Intervention
	Access}\label{practical-qualification-under-sparse-intervention-access}

In many domains, direct interventions are sparse, ethically constrained,
or expensive. This changes the evidential strategy without invalidating
the criterion. The evidential order is asymmetric: predictive success
can raise confidence, but only interventional stability can underwrite
objecthood commitment.

Where intervention data are limited, predictive closure functions as an
operational indicator while interventional closure remains the
ontological target. A model can fit historical trajectories and still
fail under novel perturbation. For that reason, claims should be
qualified by available intervention access and downgraded when
out-of-regime behavior is unstable. In-sample prediction is therefore
not treated as decisive. The standard remains robustness under
admissible perturbation, even when that robustness must be assessed
indirectly.

The same gap between in-sample fit and perturbation behavior is one
reason to triangulate diagnostics. Predictive closure can look strong
in-sample while interventional closure fails under admissible
perturbation. Partial observability can hide within-class heterogeneity
that later appears as instability. Nonstationarity can make a previously
low-leak partition unstable outside its calibration window.
Triangulation does not remove these risks, but it makes them visible.

\subsection{Conceptual Contrast: Near-Tie Prediction, Different
	Ontology
	Verdicts}\label{conceptual-contrast-near-tie-prediction-different-ontology-verdicts}

Two partitions can show similar short-horizon observational performance
and still receive different closure verdicts. Such cases show what the
criterion adds beyond predictive fit.

Suppose partition \(P_1\) and partition \(P_2\) produce comparable
one-step observational fit in calibration data. Under fixed admissible
interventions, \(P_1\) preserves stable macro-transition structure while
\(P_2\) requires repeated within-class micro refinements to maintain
performance. The two may look close observationally while differing in
transition structure.

On the present account, \(P_1\) receives the stronger objecthood
verdict because it remains transition-autonomous under the fixed
constraints. \(P_2\) can remain useful as a representation, but its
dependence on recurring hidden repair counts against macro-object
standing.

In traffic modeling, for example, one partition may track density and
flow by lane segment. A rival partition may track arbitrary vehicle-ID
clusters that happen to fit one week of observations. In-sample one-step
fit can be similar. Under admissible perturbations, such as ramp
metering changes and speed-limit adjustments, the lane-flow partition
remains stable while the ID-cluster partition becomes brittle. Under
modest horizon extension, leakiness for the ID-cluster partition rises
sharply. The verdict is then robust or qualified objecthood for the
lane-flow partition and downgrade for the rival.

Set out under the reporting schema, the case runs as follows.

\begin{enumerate}
	\item
	      Target partitions: $P_{lane}$ (lane-segment density and flow)
	      versus $P_{id}$ (vehicle-ID clusters).
	\item
	      Regime and horizon: weekday commuter traffic in one metropolitan
	      corridor; horizon set to short-run control windows and one modest
	      extension beyond calibration.
	\item
	      Admissibility class: control channels that traffic operators
	      actually use in that regime (ramp metering and speed-limit
	      controls), excluding interventions that require implausible
	      vehicle-level actuation.
	\item
	      Diagnostics: predictive gain from within-class microfeatures,
	      within-class transition-profile divergence, and intervention
	      response invariance under admissible perturbations.
	\item
	      Verdict: $P_{lane}$ remains comparatively stable and supports robust
	      or at least qualified objecthood in the declared regime; $P_{id}$
	      remains instrumentally useful in calibration but is downgraded once
	      perturbation and modest horizon extension are applied.
\end{enumerate}

The criterion also sorts between partitions that are both scientifically
rigorous. Consider two ways to divide an
organism into macroscopic parts. One partition, \(P_{\text{organ}}\),
follows functional physiological boundaries: heart, lungs, liver. A
rival partition, \(P_{\text{slice}}\), divides the body into strict
transverse geometric planes, as used in medical tomography. Under
resting observation, both achieve comparable short-horizon compression of
local mass, fluid density, and electrical activity. Under an admissible
physical perturbation, such as cardiovascular drug administration or
sudden exertion, their transition autonomy diverges.
\(P_{\text{organ}}\) preserves its macro-transition structure. The heart
responds as a coordinated unit and its macroscopic variables remain
sufficient for predicting onward causal trajectory.
\(P_{\text{slice}}\) fractures, because functional tissues span across
geometric planes; predicting the macro-future of a given slice now
requires tracking the specific intercellular signals and fluid volumes
crossing its boundary. On the present view,
\(P_{\text{organ}}\) qualifies for robust objecthood in the physiological
regime because it isolates transition-autonomous parts.
\(P_{\text{slice}}\) remains a valuable diagnostic tool, but its
recurring dependence on hidden micro-state repair prevents it from
qualifying as a macro-object under these constraints.

\subsection{Reporting and Evidential
	Discipline}\label{reporting-and-evidential-discipline}

To keep conclusions comparable across cases, closure assessments should
report five items explicitly.

\begin{enumerate}
	\item
	      Target partition and rationale.
	\item
	      Regime and horizon specification.
	\item
	      Admissible intervention class and justification.
	\item
	      Diagnostics used and their agreement or disagreement.
	\item
	      Final verdict category: robust, qualified, or indeterminate.
\end{enumerate}

Reporting these items prevents silent shifts in target, timescale, or
admissibility from being mistaken for ontological progress. The
assessment itself should follow an explicit sequence.

\begin{enumerate}
	\item
	      Specify candidate partitions for one target process.
	\item
	      Fix regime, horizon \(L\), admissible intervention class
	      \(\mathcal{I}\), diagnostics, and model classes in advance.
	\item
	      Evaluate all candidates comparatively under that same fixed setup.
	\item
	      Report verdict category and sensitivity under modest perturbation
	      checks.
\end{enumerate}

Because it fixes the setup before candidates are scored, this sequence
is part of what makes closure claims epistemically disciplined rather
than post hoc. It is not an optional implementation preference.

Diagnostics provide evidence about whether a candidate satisfies the
criterion under the stated constraints. They do not replace
criterion-level argument, and the proposal does not reduce ontology to
whichever metric happens to be available. The evidential logic is
comparative and convergent, and no single proxy is treated as
infallible. Confidence increases when different diagnostics track the
same rank-order among candidate partitions and when those rankings
remain stable under modest perturbation tests. Confidence decreases when
diagnostics diverge persistently or when rankings are brittle under
small design changes.

Verdict categories are graded accordingly. Robust verdicts require
convergence and stability, qualified verdicts fit mixed but non-trivial
evidence, and indeterminate verdicts fit persistent instability. The
diagnostic set can vary by domain, but it should answer one fixed
question: do within-class micro-differences still matter for what
happens next at the macro level? In practice, three checks are usually enough.

\begin{enumerate}
	\item
	      Leakiness proxy based on predictive gain from added within-class
	      micro-features.
	\item
	      Within-class transition-profile divergence.
	\item
	      Intervention-response invariance under admissible perturbation.
\end{enumerate}

Agreement across these checks increases warrant. Persistent disagreement
is evidence for revision or downgrade, not for forced commitment. When
interventional evidence is sparse or unavailable, the remaining
diagnostics can still support qualified status, but robust objecthood
requires that the interventional leg not be entirely missing.

As a minimal triangulation rule, a robust verdict requires stability
under modest horizon variation and under at least two defensible
admissibility perturbations, with at least two diagnostics agreeing on
candidate rank-order. If this condition fails, the verdict is qualified
or indeterminate.

\section{Objections and Replies}\label{objections-and-replies}

The objections run from foundational worries to worries about
evidential discipline. Across them the core answer stays the same:
closure under fixed constraints supplies ontological discipline without
demanding a single privileged descriptive level.

Three standing commitments apply throughout the replies below and will
not be restated each time. First, the account does not deny
microphysical completeness; it claims non-redundancy at the explanatory
and interventional level for declared macro-targets. Second, all
verdicts are indexed to the declared constraints, so objections that
presuppose context-free ontological claims address a view the paper does
not hold. Third, the strongest conclusion is conditional on structural
realist and interventionist commitments; readers who withhold those
commitments can still accept the narrower anti-gerrymandering result.

\subsection{``Instrumentalism and
	Admissibility''}\label{instrumentalism-and-admissibility}

Objection: target selection is interest-relative, so the view still
collapses into usefulness.

Reply: target selection can be interest-shaped without making closure
verdicts preference-shaped. Once regime, horizon, and admissibility are
fixed, whether macro-transitions are autonomous is a system fact under
explicit constraints.

The stronger version invokes Dennett's Martian \citep[p.~42]{dennett1991}. If an ideal
micro-predictor can forecast everything, macro realism looks optional.
The mistake is to infer ontological redundancy from representational
power. Even a micro-omniscient predictor coexists with objectively
better macro-bottlenecks for target dynamics. Ignoring those
bottlenecks increases representational burden without erasing macro
transition structure. Greater observer power changes what is convenient
to compute, but whether a partition is transition-autonomous under fixed
constraints does not vary with who computes it.

A second version targets admissibility: admissibility criteria already
encode what counts as the relevant level. The reply appeals to
hierarchical order and disclosure rules. Admissibility is fixed
by physically realizable, regime-preserving control channels before
partition scoring. Post hoc admissibility revision triggers restart.
This makes back-fitting explicit and sanctionable.

Residual disagreement can remain. The account does not promise
universal convergence from one pass. It promises disciplined comparison
and explicit downgrade when verdicts are unstable across nearby
defensible admissibility specifications. For a skeptical reader, this
means hard cases are flagged as unresolved instead of being settled by
stipulation.

Instrumentalist pressure takes two forms. One holds that ontology should
track predictive utility only; the other holds that ontology is
unnecessary once predictive utility is available. The present view
rejects both. Utility without closure is too permissive, and closure
without ontology understates what stable intervention-guiding structure
provides. What closure adds is discipline about when to commit, and that
discipline involves no metaphysical extravagance. If a candidate
repeatedly supports stable counterfactual control under fixed
constraints, refusing any ontological standing starts to look like an
empty verbal policy rather than a substantive alternative. The paper
does not force maximal realism; it requires explicit criteria for when
anti-realism remains credible.

The same reasoning answers the cheap-coding objection. A recoding can
improve local fit and still fail closure under admissible perturbation.
The criterion therefore screens out cheap recodings of this kind, which
is what permissiveness critiques demand of it.

\subsection{``Causal Exclusion Still Defeats Macro
	Claims''}\label{causal-exclusion-still-defeats-macro-claims}

Objection: if microphysics is causally complete, macro-level causal
claims are redundant \citep{kim1998}.

Reply: closure does not posit a second fundamental cause layer. It
identifies when macro-description is sufficient for prediction and
intervention at the relevant target level. Microphysical implementation
remains untouched.

The standard reply to Kim's exclusion argument is that macro-causation
is explanatorily indispensable. The closure criterion allows a stronger,
structural reply: macro-states are non-redundant because they filter out
causal degeneracy. Due to multiple realizability, myriad distinct
micro-configurations map to identical macro-outcomes. Tracking every
electron is computationally expensive, and it also introduces
non-difference-making noise into the causal model. When a system
achieves closure, the macro-boundary strips out this micro-level
degeneracy, effectively concentrating causal power. The macro-variable
is therefore a more informative causal parameter for the target outcome,
because it isolates the relevant difference-maker. It is not a
convenient shorthand competing with a fundamental micro-cause.
Supervenience establishes a dependency relation, not an identity
relation, and causal priority follows the level where conditional
independence actually screens off this noise.

This reply also handles the fat-handedness worry that
macro-interventions inevitably co-intervene on micro-states
\citep{dewhurst2021}. Co-intervention on the micro-level is
real, but closure ensures that the within-class micro-differences it
produces do not change macro-transitions; the co-intervention is
therefore non-difference-making for the macro target. Macro-level
interventions can be successful and projectible across perturbations
without denying microphysical realization \citep{woodward2003,pearl2000,hoel2013}.
The account is built to be compatible with microphysical completeness
from the outset; that compatibility is not a concession.

Nor does the account face an overdetermination problem. It claims one
implemented process with multiple adequate descriptions at different
explanatory and control targets, not two independent sufficient causes
for one event. Exclusion pressure weakens once adequacy is indexed to
target and intervention class instead of to a single privileged level.

Against Kim's strongest formulations, the account indexes adequacy to
targets instead of multiplying layers \citep{kim1998,kim2005}. A
closed macro-description can be non-redundant without competing for
micro-level fundamentality, because non-redundancy is inferred when
micro-detail ceases to add control-relevant information for the target
outcome under fixed constraints. The claim concerns which
counterfactual dependencies remain stable for a declared class of
interventions, so it is not a linguistic escape. Compatibility with
microphysical completeness by itself is too weak to do this work,
because it leaves open whether macro-variables are dispensable
shorthand. Closure narrows that space. When it holds, dispensability
must be argued against a stated record of interventional and predictive
sufficiency, not simply asserted.

\subsection{``Autonomy, Pluralism, and
	Distinctiveness''}\label{autonomy-pluralism-and-distinctiveness}

Objection: the view sounds like a formal restatement of familiar
special-sciences autonomy claims, not a distinctive realism thesis.

Reply: the paper is continuous with autonomy insights, but it is not
equivalent to them. Generic autonomy claims often remain permissive
about what qualifies as a level-worthy grouping. The present view adds a
discriminating closure condition with explicit exclusion and downgrade
rules, an addition that changes verdict structure.

The difference is both practical and philosophical. The account gives a
protocol for ruling out high-maintenance, transition-incoherent
candidates, and it ties realism commitment to transition autonomy under
admissible interventions rather than to explanatory convenience alone.
Causal-emergence style results are accordingly treated as evidence
within a closure-governed approach, not as a replacement for the
criterion itself \citep{hoel2013}.

Mechanistic accounts are right to tie level talk to organized structure
that supports stable intervention and explanation
\citep{machamer2000,craver2007}. They do not, by themselves, settle
which variables or entities have been carved at the right grain for the
target explanation. Closure addresses that earlier question by testing
whether a candidate macro-variable screens off the micro-differences
that would otherwise keep the mechanism explanatorily open.

The view is therefore best read as a constrained realism refinement of
autonomy views, neither an unrelated alternative nor a mere
restatement. Generic autonomy claims assert that higher levels matter.
Closure specifies \emph{when} they matter, \emph{how much} they matter
(graded by leakiness), and \emph{when the claim should be withdrawn}
(explicit downgrade conditions). Without withdrawal conditions, a
proposed realism criterion could affirm commitment but could never
count against a candidate.

The classical arguments for the autonomy of the special sciences
established that higher-level kinds can be explanatorily indispensable
because they abstract over heterogeneous micro-realizations
\citep{putnam1967,fodor1974}. Multiple realizability is in fact a
structural feature of any closed macro-kind: if a partition closes,
different microstates within the same macroclass by construction
produce the same macro-transitions, making the kind multiply realized
in the relevant sense. Multiple realizability is therefore necessary
for closure, but it is not sufficient. A purely extensional claim about
kind membership says nothing about the dynamical coherence of the
realizers. A disjunctive aggregate can have heterogeneous realizers
while completely failing closure, because those realizers can have
utterly divergent transition profiles. Closure adds a dynamical
requirement. Falling under the same label is not enough; the realizers
must produce the same macro-transitions under the stated constraints.
Gerrymandered kinds exploit multiple-realizability intuitions in the gap
between what a kind groups together extensionally and how those
instances behave dynamically, without earning macro-object status. The present criterion formalizes the transition condition that
the autonomy tradition always needed to separate genuine higher-level
kinds from merely disjunctive aggregates \citep{shapiro2000}. The
monetary illustration in \S\ref{one-distributed-illustration}
provides a concrete instance. Coins, ledger entries, and digital
balances are heterogeneous realizers, but money's macro-object status
rests on the transition-equivalence of those realizers under admissible
financial interventions, and their heterogeneity alone would not
secure it.

Pluralism version of the objection: if multiple grains can satisfy
closure, ontology becomes permissive again. Reply: plurality of closed
grains does not imply arbitrariness. Once regime, horizon, and
admissibility are fixed, which partitions close is determined by
transition structure, not by analyst choice. When multiple candidates
remain viable, robustness and minimality rank them as an objective
relation among candidates under fixed constraints. In many domains these
relations are further structured by coarse-graining order. When one
closed partition is a refinement of another, the two nest instead of
competing, which gives pluralism structure without dissolving it.

Plurality also holds within a single physical system. The same substrate
can support several closed coarse-grainings in overlapping regimes, and
each can meet the objecthood condition in its own regime. Ordinary kind
terms need not track these structures one to one. As a conceptual
illustration, suppose a lump of wood and metal formed by accident and,
under the interventions involved in striking and pulling nails, behaves
exactly as a manufactured hammer does. In that regime it realizes the
same transition structure as the manufactured hammer, and the criterion
treats the two alike. Whether the lump is a hammer in the artifact
sense, given that nobody made it for that purpose, is a separate
question. Ordinary language bundles the two questions under one word,
and the criterion answers only the dynamical one. The artifact question
is left to the accounts of classification and individuation noted in
\S\ref{scope-and-non-claims}.

Admitting this plurality suits complex systems. Forcing a single grain
in every context would be a stronger and less defensible claim than the
paper needs.

\subsection{``The Markov Template is Too
	Restrictive''}\label{the-markov-template-is-too-restrictive}

Objection: strong lumpability presupposes first-order Markov
microdynamics and excludes memory-dependent systems.

Reply: the criterion is transition sufficiency, not first-order
Markovity. Strong lumpability is an exact benchmark. In memory-bearing
systems, state can be enriched with relevant history and the same
closure question can be asked over that enriched state. Only the state
representation changes; the content of the criterion stays the same.

Immune response modeling, for example, often depends on historical
exposure profiles, and institutional dynamics often depend on
path-dependent enforcement and trust trajectories. In both cases, a
memoryless state is too thin, but a minimally history-enriched state can
still support closure testing at the target grain.

Minimality still matters. Enrichment should be the least extension
needed to preserve closure performance under fixed constraints. The
account also accepts an edge case: if the minimally sufficient enriched
state approaches micro-level complexity, then closure may be recovered
without meaningful compression. In that regime the result does not
support robust macro-objecthood. It signals instead that no informative
macro-partition has been found at the target grain.

\subsection{``Formal Closure Collapses the View into
	Formalism''}\label{formal-closure-collapses-the-view-into-formalism}

Objection: formal systems are closed by stipulation, so closure cannot
ground empirical objecthood.

Reply: stipulated and induced closure must be separated. Formal systems
can have objective internal closure under constitutive rules. The
present criterion concerns induced closure in implemented spatiotemporal
dynamics. Internal formal coherence does not by itself settle empirical
macro-objecthood claims.

This does not reduce formal systems to arbitrary invention. Formal
systems are invented, but not freely. What counts as intelligible formal
practice is constrained by stable identity conditions, compositional
inference, and coherent consequence. Those constraints explain why
formal work can guide empirical reasoning without itself deciding
empirical objecthood.

The view therefore does not collapse into formalism. Its claim is the
narrower one that induced closure is the objecthood criterion for
implemented macro-patterns in regime.

\subsection{``Closure Is Too Strong and Eliminates Most Special
	Sciences''}\label{closure-is-too-strong-and-eliminates-most-special-sciences}

Objection: if closure requires that within-class micro-differences not
matter for macro-transitions, then most special-science kinds will fail.
Biological species, psychological states, and economic categories are
notoriously leaky. The criterion eliminates the ontology it was supposed
to discipline.

Reply: this objection conflates robust objecthood with any objecthood
verdict at all. The account provides three verdict categories, not
one. Robust objecthood requires stable closure under fixed constraints.
Qualified objecthood fits cases where closure is partial,
regime-limited, or supported by convergent but incomplete evidence.
Indeterminate status fits cases where evidence is too unstable to
warrant commitment.

Most special-science kinds fall into the qualified range, and that is
the correct result. Biological species exhibit substantial transition
autonomy within ecological and evolutionary regimes while leaking at
boundary cases (hybrid zones, ring species, horizontal gene transfer).
Psychological kinds often support predictive and interventional
regularity within clinical or experimental regimes while failing under
broader perturbation. The criterion does not eliminate these kinds. It
assigns them the status their closure evidence supports, which is
qualified objecthood with explicit regime limitations, rather than either full robust status or wholesale elimination.

The objection therefore proves too much. A criterion that granted robust
objecthood to every special-science kind regardless of leakiness would
be the permissive view this paper is designed to replace. Graded
verdicts can match the evidential situation instead of forcing a binary
choice between full realism and eliminativism.

\subsection{``Predictive Success and Failure
	Conditions''}\label{predictive-success-and-failure-conditions}

Objection: a model can predict well without warranting ontological
commitment.

Reply: the account accepts this. Predictive success functions as
evidence only when it aligns with closure diagnostics and
interventional stability. Two consequences follow.

First, protocol language has a bounded philosophical role. The paper
includes diagnostics to avoid a familiar failure mode, where a criterion
is asserted but left too unconstrained to guide verdicts. The claim
itself remains conceptual. The diagnostics operationalize what the
argument says should matter; they do not replace the argument.

Second, the criterion is falsifiable in its own terms. If candidate
partitions repeatedly fail closure diagnostics across defensible regimes
and admissibility classes, commitment should be withdrawn or downgraded.
The same applies when verdicts are unstable under small, defensible
changes in horizon, admissibility, or diagnostic proxy. A proposal that
could not risk downgrade would not be a serious realism criterion.

The proposal offers explicit criteria for when confidence increases,
when it should pause, and when commitment should retreat. It does not
promise certainty from one metric or one pass.

\subsection{``Horizon-Relativity Makes the Criterion Too
	Weak''}\label{horizon-relativity-makes-the-criterion-too-weak}

Objection: if closure is indexed to horizon, any candidate can be made
to pass at a short enough horizon.

Reply: two requirements keep horizon indexing from becoming an escape
clause.

First, horizons must be anchored in independently identified timescale
structure in the target domain: known relaxation windows, intervention
latency, control-response periods, or similar physically motivated
temporal features. A horizon chosen only because it makes a favored
partition look closed has no independent standing.

Second, verdicts must survive a modest-extension stability check. A
candidate that passes at horizon $L$ but fails immediately at $L +
	\Delta$ for small, defensible $\Delta$ receives, at best, qualified
status. Robust objecthood requires that closure not be razor-thin in
temporal scope.

Together, these requirements force explicit timescale discipline and
prevent silent switching between incompatible temporal targets. Horizon
flexibility is acceptable when it tracks independently motivated
structure. It is not acceptable when it is introduced only after score
inspection to salvage a preferred partition.

Critics are right that any criterion can be trivialized if timescales
are unconstrained. The present account avoids that outcome by tying
horizons to ex ante justification and by requiring modest-extension
robustness. If robustness fails, the verdict is downgraded, and the
criterion has then done its job by recording that the claim was too
strong for the evidence.

\subsection{``The Criterion Is Too
	Conservative''}\label{the-criterion-is-too-conservative}

Objection: by demanding closure and downgrade discipline, the criterion
may be too conservative and may miss emerging macro-objects.

Reply: conservatism is partly intentional. The paper aims to avoid
premature ontological inflation. Still, the criterion is not rigid. It
allows qualified verdicts for emerging patterns when closure evidence is
partial but improving, and it allows verdict revision as regimes
stabilize and intervention evidence accumulates.

So the criterion does not require all-or-nothing maturity before any
commitment. It requires that commitment strength track available closure
evidence. \citet{woodward2021} similarly argues that while minimal
interventionist criteria establish causality, invariance is a graded
virtue. Our criterion of closure can be seen as the dynamical
realization of this: a macro-object is ``proportional'' when it
minimizes the omission of relevant micro-detail. This graded structure
suits dynamic systems, where objecthood can be developmental rather
than instantaneous.

Compression and closure also reconnect here in a non-vacuous way. Early
in a domain's development, compression gains may appear before robust
closure is demonstrated. The account treats this as evidential lead,
not ontological conclusion. As closure evidence accumulates, commitment
can strengthen. If closure evidence fails to accumulate, commitment
should remain qualified or be withdrawn even when short-run compression
remains attractive. The criterion is therefore neither eliminativist nor
inflationary. It does not deny standing to emerging patterns, and it
does not award standing before the dynamical evidence warrants it.

\section{Conclusion: A Disciplined Realism
  Claim}\label{conclusion-a-disciplined-realism-claim}

Real-pattern realism is plausible but underconstrained. This paper adds
the missing discriminating condition: closure of macro-transition
structure under fixed regime, horizon, and admissible intervention
class. Strong lumpability provides the exact benchmark. Leakiness and
convergent diagnostics extend the same criterion to non-ideal settings
without changing its content.

The proposal supports selective commitment: robust commitment where
closure is stable, qualified commitment in borderline cases, and
explicit downgrade where verdicts depend on fragile admissibility or
horizon choices. It is intentionally middle-range. It is stronger than
compression-only pattern realism because it adds exclusion and
downgrade structure, and weaker than a total-level metaphysics because
it is restricted to induced closure in spatiotemporal systems under
declared constraints.

The view also resists eliminativist pressure. It does not treat
higher-level descriptions as fictions by default, and it does not grant
them standing by convenience alone. It asks whether they track stable
transition structure under admissible intervention.

Clear limits remain. The proposal is conditional on the structural
realist and interventionist commitments stated in
\S\ref{introduction-why-pattern-realism-needs-a-stricter-criterion};
readers who withhold those commitments can still accept the narrower
anti-gerrymandering result. The objecthood it secures is the existence
of an autonomous macro-level unit; questions of kind, origin, function,
and individuation are left to the complementary accounts noted in
\S\ref{scope-and-non-claims}. It is not a complete metaphysics of
levels and not a universal methods manual. Its contribution is a
philosophical criterion with application discipline. Methodological
implementation details remain downstream of this argument, and refining
them does not alter the criterion's philosophical content. Future work should
test the criterion across domains by comparing candidate partitions
under fixed constraints, rather than by multiplying new concepts. If
rejected, it should be rejected because one rejects its stated
commitments, not because circularity, instrumentalist drift, or scope
ambiguity were left unresolved.

\backmatter

\label{references}
